\section{Motivating Example}
\label{sec:background}

\Cref{fig:app} (\cref{sec:eval}) shows \texttt{broken-admin}, a six-page
application built as ground truth. A guest lands on \texttt{index}, signs in,
and reaches \texttt{dashboard}. From there the guest can open
\texttt{reports} and \texttt{admin}, and every page links back. The policy,
written in English and stored with its structured form, reads:

\begin{quote}\small
(F1) A guest must never see the admin link.
(F2) A guest must never be able to delete a user.
(F3) No page may ship a debug panel.
(R1) Every page must offer a skip-to-content link.
\end{quote}

Four violations are planted. The Admin link on \texttt{dashboard} breaks F1.
The debug panel on \texttt{reports} breaks F3. The Delete-user button on
\texttt{admin} breaks F2. \texttt{admin} is also the only page without a
skip-to-content link, which breaks R1.

Three observations follow from this example. \emph{First}, no workflow test
catches F1--F3, because each defect is extra UI, not broken UI. \emph{Second},
R1 is an \emph{absence}. An engine that only intersects what it sees with a
forbidden list cannot express it (\cref{sec:eval:rq3}). \emph{Third}, the
graph is dense with cycles (log out, back to dashboard). A crawler that drops
back-edges still ``finds the violations'' but throws away the navigation
evidence that a path-sensitive policy (``reachable without signing in'') would
need.
